\documentclass{Phyport}

\usepackage{hyperref}
\hypersetup{
    colorlinks=true,
}
% 选用宏（如有额外需要的宏包请在此添加）
\usepackage{float}       % 提供 [H] 选项
\usepackage{booktabs}    % 三线表
\usepackage{amsmath}     % 数学公式
\usepackage{caption}     % 标题设置
 %选用宏

\exname{用霍尔法测直流圆线圈与亥姆霍兹线圈磁场} %实验名称
\instructor{周小舟} %指导教师
\class{电科2502} %班级
\name{管云哲} %姓名
\stuid{3250104519} %学号

\nyear{2026} %年
\nmonth{9} %月
\nday{20} %日
\nweekday{天} %星期几

\redate{} %如有实验补做，补做日期
\resitu{} %情况说明：

\begin{document}
\setcounter{page}{0}
\makecover

% ============================================================
%                  一、预习报告（10分）
% ============================================================
\section{预习报告（10分）}
（注：将已经写好的"物理实验预习报告"内容拷贝过来）

\subsection{实验综述（5分）}
（自述实验现象、实验原理和实验方法，包括必要的光路图、电路图、公式等。不超过500字。）

实验现象：磁场强度随电流强度与空间位置变化，亥姆霍兹线圈中部存在均匀磁场区域。

实验原理：在$x$方向通电流$I_S$，在$z$方向加磁场$B$，则在$y$方向，即试样P、S电极两侧聚积等量异号电荷，产生附加电场。电子受洛伦兹力$f_B=evB$，受侧向电场力$f_E=eE_H$。稳态时两力平衡，可得
\begin{equation}
    V_H = R_H \frac{I_S B}{d} = K_H I_S B
\end{equation}
其中$R_H=1/(ne)$为霍尔系数，$K_H$为霍尔灵敏度。由此得磁场测量公式：
\begin{equation}
    B = \frac{V_H}{K_H I_S}
\end{equation}
亥姆霍兹线圈由两个间距等于半径$\bar{R}$的同轴圆线圈组成，轴线上的磁场为
\begin{equation}
    B = \frac{1}{2}\mu_0 N I \bar{R}^2 \left[\left(\bar{R}^2 + \left(\frac{\bar{R}}{2}+x\right)^2\right)^{-3/2} + \left(\bar{R}^2 + \left(\frac{\bar{R}}{2}-x\right)^2\right)^{-3/2}\right]
\end{equation}
其中点处$B_0 = \left(\dfrac{4}{5}\right)^{3/2} \dfrac{\mu_0 N I}{\bar{R}}$。

实验方法：1、测绘单个圆线圈磁场强度分布；2、测绘亥姆霍兹线圈磁场强度分布。

\subsection{实验重点（3分）}
（简述本实验的学习重点，不超过100字。）

理解霍尔效应基本原理及磁场测量方法，掌握利用霍尔传感器测量单个圆线圈和亥姆霍兹线圈磁场分布的实验技术，学会绘制与分析$B$-$x$曲线，验证亥姆霍兹线圈中部的均匀磁场特性。

\subsection{实验难点（2分）}
（简述本实验的实现难点，不超过100字。）

精确调节霍尔传感器的位置与方向，确保其与磁场方向垂直；排除地磁场等外部磁场干扰；正确处理霍尔电压与磁场强度的关系，验证理论公式与实验结果的吻合度。

% ============================================================
%                  二、原始数据（20分）
% ============================================================
\section{原始数据（20分）}
（将有老师签名的"自备数据记录草稿纸"的扫描或手机拍摄图粘贴在下方，完整保留姓名，学号，教师签字和日期。）

\begin{figure}[H]
    \centering
    \includegraphics[width=0.85\linewidth]{实验数据}
    \caption{原始数据}
    \label{fig:rawdata}
\end{figure}

% ============================================================
%                  三、结果与分析（60分）
% ============================================================
\section{结果与分析（60分）}

\subsection{数据处理与结果（30分）}
（列出数据表格、选择适合的数据处理方法、写出测量或计算结果。）

\subsubsection{实验参数}

本实验所用参数如表~\ref{tab:params} 所示。

\begin{table}[H]
    \centering
    \caption{实验参数}
    \label{tab:params}
    \begin{tabular}{cccc}
        \toprule
        参数 & 符号 & 数值 & 单位 \\
        \midrule
        线圈匝数 & $N$ & 400 & 匝 \\
        励磁电流 & $I$ & 0.400 & A \\
        线圈半径 & $R$ & 0.100 & m \\
        真空磁导率 & $\mu_0$ & $4\pi \times 10^{-7}$ & T$\cdot$m/A \\
        \bottomrule
    \end{tabular}
\end{table}

\subsubsection{地磁场测量}

测得地磁场磁感应强度 $B_{\text{earth}} = 0.028\;\text{mT} = 28\;\mu\text{T}$。

为消除地磁场等系统误差的影响，每组测量均采用正向励磁电流与反向励磁电流分别测量，取两次测量的绝对值之平均作为最终结果：
\begin{equation*}
    B_{\text{平均}} = \frac{|B_{\text{正}}| + |B_{\text{反}}|}{2}
\end{equation*}

\subsubsection{载流圆线圈轴向磁场分布}

测量条件：励磁电流 $I = 0.400\;\text{A}$，霍尔探头沿线圈轴线移动。测量数据及与理论值的对比见表~\ref{tab:axial}，$B$-$X$ 曲线见图~\ref{fig:axial}。

单线圈轴线上磁感应强度的理论公式为
\begin{equation}
    B(x) = \frac{\mu_0 N I R^2}{2(R^2 + x^2)^{3/2}}
    \label{eq:single}
\end{equation}

\begin{table}[H]
    \centering
    \caption{载流圆线圈轴向磁场分布}
    \label{tab:axial}
    \begin{tabular}{ccccc}
        \toprule
        $X$ / cm & $B_{\text{正}}$ / mT & $B_{\text{反}}$ / mT & $B_{\text{平均}}$ / mT & $B_{\text{理论}}$ / mT \\
        \midrule
        $-$5.00 & 0.720 & $-$0.720 & 0.720 & 0.719 \\
        $-$4.00 & 0.804 & $-$0.806 & 0.805 & 0.805 \\
        $-$3.00 & 0.882 & $-$0.884 & 0.883 & 0.883 \\
        $-$2.00 & 0.942 & $-$0.944 & 0.943 & 0.948 \\
        $-$1.00 & 0.988 & $-$0.988 & 0.988 & 0.990 \\
         0.00   & 1.006 & $-$1.002 & 1.004 & 1.005 \\
         1.00   & 0.990 & $-$0.988 & 0.989 & 0.990 \\
         2.00   & 0.946 & $-$0.946 & 0.946 & 0.948 \\
         3.00   & 0.884 & $-$0.884 & 0.884 & 0.883 \\
         4.00   & 0.808 & $-$0.804 & 0.806 & 0.805 \\
         5.00   & 0.724 & $-$0.720 & 0.722 & 0.719 \\
        \bottomrule
    \end{tabular}
\end{table}

\begin{figure}[H]
    \centering
    \includegraphics[width=0.75\linewidth]{single_coil_axial}
    \caption{载流圆线圈轴向磁场分布（红点为实验值，蓝线为理论曲线）}
    \label{fig:axial}
\end{figure}

由表~\ref{tab:axial} 和图~\ref{fig:axial} 可知，实验值与理论值吻合良好，线圈中心（$X = 0$）处磁感应强度最大，$B_{\text{平均}} = 1.004\;\text{mT}$，数据关于 $X = 0$ 近似对称，远离中心处磁场衰减，符合理论预期。

\subsubsection{载流圆线圈径向磁场分布}

测量条件：霍尔探头沿线圈径向（垂直于轴线方向）移动。测量数据见表~\ref{tab:radial}，$B$-$Y$ 曲线见图~\ref{fig:radial}。

\begin{table}[H]
    \centering
    \caption{载流圆线圈径向磁场分布}
    \label{tab:radial}
    \begin{tabular}{cccc}
        \toprule
        $Y$ / cm & $B_{\text{正}}$ / mT & $B_{\text{反}}$ / mT & $B_{\text{平均}}$ / mT \\
        \midrule
        $-$5.00 & 1.240 & $-$1.236 & 1.238 \\
        $-$4.00 & 1.138 & $-$1.132 & 1.135 \\
        $-$3.00 & 1.078 & $-$1.068 & 1.073 \\
        $-$2.00 & 1.036 & $-$1.028 & 1.032 \\
        $-$1.00 & 1.012 & $-$1.006 & 1.009 \\
         0.00   & 1.006 & $-$1.000 & 1.003 \\
         1.00   & 1.016 & $-$1.006 & 1.011 \\
         2.00   & 1.042 & $-$1.028 & 1.035 \\
         3.00   & 1.084 & $-$1.076 & 1.080 \\
         4.00   & 1.148 & $-$1.140 & 1.144 \\
         5.00   & 1.256 & $-$1.244 & 1.250 \\
        \bottomrule
    \end{tabular}
\end{table}

\begin{figure}[H]
    \centering
    \includegraphics[width=0.75\linewidth]{single_coil_radial}
    \caption{载流圆线圈径向磁场分布}
    \label{fig:radial}
\end{figure}

由表~\ref{tab:radial} 和图~\ref{fig:radial} 可知，线圈中心（$Y = 0$）处磁感应强度最小（$B_{\text{平均}} = 1.003\;\text{mT}$），径向磁场随 $|Y|$ 增大而增大，这是因为远离轴线处磁力线弯曲导致径向分量增大。数据关于 $Y = 0$ 近似对称。

\subsubsection{亥姆霍兹线圈轴向磁场分布}

测量条件：两线圈间距 $d = R = 0.100\;\text{m}$，构成亥姆霍兹线圈。霍尔探头沿公共轴线移动，测量范围 $X = -12\;\text{cm}$ 至 $+12\;\text{cm}$。测量数据及与理论值的对比见表~\ref{tab:helmholtz}，$B$-$X$ 曲线见图~\ref{fig:helmholtz}。

亥姆霍兹线圈轴线上磁感应强度的理论公式为
\begin{equation}
    B(x) = \frac{\mu_0 N I R^2}{2(R^2 + (x + R/2)^2)^{3/2}} + \frac{\mu_0 N I R^2}{2(R^2 + (x - R/2)^2)^{3/2}}
    \label{eq:helmholtz}
\end{equation}

\begin{table}[H]
    \centering
    \caption{亥姆霍兹线圈轴向磁场分布}
    \label{tab:helmholtz}
    \small
    \begin{tabular}{ccccc}
        \toprule
        $X$ / cm & $B_{\text{正}}$ / mT & $B_{\text{反}}$ / mT & $B_{\text{平均}}$ / mT & $B_{\text{理论}}$ / mT \\
        \midrule
        $-$12.00 & 0.682 & $-$0.670 & 0.676 & 0.684 \\
        $-$11.00 & 0.782 & $-$0.776 & 0.779 & 0.784 \\
        $-$10.00 & 0.886 & $-$0.884 & 0.885 & 0.891 \\
         $-$9.00 & 1.006 & $-$1.002 & 1.004 & 1.002 \\
         $-$8.00 & 1.116 & $-$1.110 & 1.113 & 1.111 \\
         $-$7.00 & 1.222 & $-$1.218 & 1.220 & 1.212 \\
         $-$6.00 & 1.308 & $-$1.302 & 1.305 & 1.296 \\
         $-$5.00 & 1.374 & $-$1.372 & 1.373 & 1.361 \\
         $-$4.00 & 1.422 & $-$1.420 & 1.421 & 1.403 \\
         $-$3.00 & 1.446 & $-$1.440 & 1.443 & 1.427 \\
         $-$2.00 & 1.460 & $-$1.458 & 1.459 & 1.436 \\
         $-$1.00 & 1.462 & $-$1.460 & 1.461 & 1.439 \\
          0.00   & 1.464 & $-$1.464 & 1.464 & 1.439 \\
          1.00   & 1.462 & $-$1.458 & 1.460 & 1.439 \\
          2.00   & 1.458 & $-$1.454 & 1.456 & 1.436 \\
          3.00   & 1.448 & $-$1.440 & 1.444 & 1.427 \\
          4.00   & 1.418 & $-$1.412 & 1.415 & 1.403 \\
          5.00   & 1.372 & $-$1.364 & 1.368 & 1.361 \\
          6.00   & 1.308 & $-$1.302 & 1.305 & 1.296 \\
          7.00   & 1.218 & $-$1.216 & 1.217 & 1.212 \\
          8.00   & 1.122 & $-$1.116 & 1.119 & 1.111 \\
          9.00   & 1.004 & $-$1.000 & 1.002 & 1.002 \\
         10.00   & 0.898 & $-$0.890 & 0.894 & 0.891 \\
         11.00   & 0.788 & $-$0.776 & 0.782 & 0.784 \\
         12.00   & 0.690 & $-$0.680 & 0.685 & 0.684 \\
        \bottomrule
    \end{tabular}
\end{table}

\begin{figure}[H]
    \centering
    \includegraphics[width=0.85\linewidth]{helmholtz_coil}
    \caption{亥姆霍兹线圈轴向磁场分布（红点为实验值，蓝线为理论曲线，灰色虚线标示线圈位置）}
    \label{fig:helmholtz}
\end{figure}

由表~\ref{tab:helmholtz} 和图~\ref{fig:helmholtz} 可知，两线圈之间存在明显的均匀磁场区域。在 $|X| \leq 2\;\text{cm}$ 范围内，$B$ 值变化极小（$1.456 \sim 1.464\;\text{mT}$），中心处 $B_{\text{平均}} = 1.464\;\text{mT}$ 为最大值。整体呈中间高、两侧低的单峰结构，远离中心处磁场衰减。

\subsubsection{三图对比}

为便于对比，将三组磁场分布曲线合并展示于图~\ref{fig:combined}。

\begin{figure}[H]
    \centering
    \includegraphics[width=\linewidth]{B_X_curves_combined}
    \caption{磁场分布对比：(a) 单线圈轴向；(b) 单线圈径向；(c) 亥姆霍兹线圈轴向}
    \label{fig:combined}
\end{figure}

% ============================================================
%                  误差分析（20分）
% ============================================================
\subsection{误差分析（20分）}
（运用测量误差、相对误差或不确定度等分析实验结果，写出完整的结果表达式，并分析误差原因。）

\subsubsection{载流圆线圈轴向磁场误差}

单线圈中心处（$X = 0$）的测量结果与理论值对比：
\begin{equation*}
    B_{\text{实验}} = 1.004\;\text{mT}, \quad B_{\text{理论}} = 1.005\;\text{mT}
\end{equation*}
\begin{equation*}
    \delta_r = \frac{B_{\text{实验}} - B_{\text{理论}}}{B_{\text{理论}}} \times 100\% = \frac{1.004 - 1.005}{1.005} \times 100\% = -0.1\%
\end{equation*}

各测量点的相对误差：平均绝对相对误差为 $0.2\%$，最大相对误差为 $0.5\%$（$X = -2.0\;\text{cm}$ 处）。

测量结果表达式：
\begin{equation*}
    B_0 = (1.004 \pm 0.005)\;\text{mT} \quad (\delta_r = -0.1\%)
\end{equation*}

\subsubsection{亥姆霍兹线圈轴向磁场误差}

亥姆霍兹线圈中心处（$X = 0$）的测量结果与理论值对比：
\begin{equation*}
    B_{\text{实验}} = 1.464\;\text{mT}, \quad B_{\text{理论}} = 1.439\;\text{mT}
\end{equation*}
\begin{equation*}
    \delta_r = \frac{1.464 - 1.439}{1.439} \times 100\% = 1.8\%
\end{equation*}

各测量点的相对误差：平均绝对相对误差为 $0.8\%$，最大相对误差为 $1.8\%$（$X = 0.0\;\text{cm}$ 处）。

均匀区域（$|X| \leq 3\;\text{cm}$）的实验平均值为 $1.455\;\text{mT}$，与中心理论值的相对误差为 $1.2\%$。

测量结果表达式：
\begin{equation*}
    B_0 = (1.464 \pm 0.026)\;\text{mT} \quad (\delta_r = 1.8\%)
\end{equation*}

\subsubsection{误差原因分析}

\begin{enumerate}
    \item \textbf{系统误差}：亥姆霍兹线圈的实验值系统性地高于理论值（约 $1\% \sim 2\%$），可能原因包括：
    \begin{enumerate}
        \item 霍尔传感器的灵敏度参数 $K_H$ 存在校准偏差；
        \item 工作电流 $I_S$ 的实际值与标称值存在差异；
        \item 线圈实际半径与标称值 $R = 0.100\;\text{m}$ 有微小偏差。
    \end{enumerate}
    \item \textbf{随机误差}：正向与反向电流测量值之间存在微小差异（地磁场残余影响未完全消除），以及霍尔探头定位精度有限。
    \item \textbf{单线圈误差较小}的原因：单线圈中心处仅由一个线圈贡献磁场，对参数偏差不敏感；而亥姆霍兹线圈为两线圈叠加，参数偏差被放大。
\end{enumerate}

% ============================================================
%                  实验探讨（10分）
% ============================================================
\subsection{实验探讨（10分）}
（对实验内容、现象和过程的小结，不超过100字。）

本实验成功利用霍尔效应测量了单线圈与亥姆霍兹线圈的磁场分布，实验值与理论吻合良好（误差$\leq$2\%），验证了亥姆霍兹线圈均匀区特性。正反电流平均法有效消除了地磁场干扰。

% ============================================================
%                  四、思考题（10分）
% ============================================================
\section{思考题（10分）}
（解答教材或讲义或老师布置的思考题，请先写题干，再作答。）

\textbf{1. 亥姆霍兹线圈是怎样定义的？什么条件下能产生均匀磁场？均匀区范围有多大？}

亥姆霍兹线圈由两个半径为$R$、匝数为$N$、通同向等大电流$I$的共轴圆线圈组成，两线圈间距等于线圈半径，即$d = R$。

当两线圈间距恰好等于线圈半径$d = R$时，轴线上磁感应强度$B(x)$在中心点$x = 0$处的二阶导数$\dfrac{\mathrm{d}^2 B}{\mathrm{d}x^2}\bigg|_{x=0} = 0$，磁场在中心附近最为均匀。这是因为对$B(x)$做泰勒展开时，常数项和二次项共同决定了均匀程度，当$d = R$时二次项恰好为零。

均匀区范围约为线圈中心两侧各$R/2$以内。本实验中$R = 0.100\;\text{m}$，均匀区约为$|X| \leq 3\;\text{cm}$（即$|X|/R \leq 0.3$），该区域内$B$的相对变化不超过$1.5\%$。

\vspace{1em}

\textbf{2. 如果霍尔传感器与磁场方向不垂直，对测量结果有什么影响？}

霍尔效应的原理是在垂直于磁场和电流的方向上产生霍尔电势差。设磁场$B$与霍尔传感器法线方向之间的夹角为$\theta$，则霍尔传感器测得的磁感应强度为有效分量：
\begin{equation*}
    B_{\text{测}} = B \cos\theta
\end{equation*}
当$\theta = 0$（传感器与磁场垂直）时，$B_{\text{测}} = B$，测量值准确。当$\theta \neq 0$时，测量值偏小，引入的系统误差为
\begin{equation*}
    \delta = \frac{B_{\text{测}} - B}{B} = \cos\theta - 1 \approx -\frac{\theta^2}{2}
\end{equation*}
对于小角度，误差与$\theta^2$成正比。例如$\theta = 5°$时，$\delta \approx -0.4\%$；$\theta = 10°$时，$\delta \approx -1.5\%$。因此安装传感器时必须确保其法线方向与被测磁场方向严格一致，否则将引入负向系统误差，使测量值偏低。

{\fangsong 注意事项：
\begin{enumerate}[label=\arabic*.]
    \item 用PDF格式上传“实验报告”，文件名：学生姓名+学号+实验名称+周次。
    \item “实验报告”必须递交在“学在浙大”本课程内对应实验项目的“作业”模块内。
    \item “实验报告”成绩必须在“浙江大学物理实验教学中心网站”-“选课系统”内查询。
    \item 教学评价必须在“浙江大学物理实验教学中心网站”-“选课系统”内进行，学生必须进行教学评价，才能看到实验报告成绩，教学评价须在本次实验结束后3天内进行。
\end{enumerate}}
\begin{center}
    \fbfangsong{浙江大学物理实验教学中心制}
\end{center}
\end{document}
